\section{Iterative Postorder Traversal (Two Stacks)}
\label{sec:trees-iter-post-two}

\problemstatement{Return the postorder traversal (left, right, node) of a
binary tree without recursion, using two stacks.}

\begin{patternbox}
Postorder is the hardest to make iterative, because a node is recorded
only \emph{after both} subtrees. The two-stack method sidesteps the
difficulty with a slick observation: postorder is the exact
\textbf{reverse} of a slightly modified preorder. If you can produce
``node, right, left'' and then reverse it, you have postorder for free.
\end{patternbox}

\subsection*{Understanding the problem carefully}

Postorder is left, right, node. Read it backwards and you get node, right,
left -- which is just preorder with the two child visits swapped. So if we
run a preorder-like traversal that pushes \emph{left before right} (giving
node, right, left order) and collect the results, then reverse the whole
collection, we recover true postorder. The second stack is simply the
container we reverse.

\begin{ideabox}[Postorder = reverse of (node, right, left)]
Run a preorder variant using stack~1 that emits nodes in \emph{node,
right, left} order, pushing each emitted node onto stack~2. Popping
stack~2 at the end yields those nodes in reverse -- i.e.\ \emph{left,
right, node}, exactly postorder. No delayed-visit bookkeeping is needed.
\end{ideabox}

\subsection*{Algorithm}

\begin{enumerate}
  \item If root is null, return empty. Push root onto stack~1.
  \item While stack~1 is non-empty:
  \begin{enumerate}
    \item Pop a node from stack~1 and push it onto stack~2.
    \item Push that node's \emph{left} child, then its \emph{right}
          child, onto stack~1 (if they exist).
  \end{enumerate}
  \item Pop everything off stack~2 into the output; that is the
        postorder.
\end{enumerate}

\subsection*{Dry run}

\dryrun{Tree: root $1$; children $2$ (left), $3$ (right).}

\begin{itemize}
  \item Stack1 $=[1]$: pop $1\to$ stack2 $=[1]$; push $2$ then $3$. Stack1
        $=[2,3]$.
  \item Pop $3\to$ stack2 $=[1,3]$ (no children). Stack1 $=[2]$.
  \item Pop $2\to$ stack2 $=[1,3,2]$. Stack1 empty.
  \item Empty stack2: $2,3,1$. That is the postorder.
\end{itemize}

\subsection*{C++ implementation}

\begin{lstlisting}
#include <bits/stdc++.h>
using namespace std;

struct TreeNode {
    int val;
    TreeNode* left;
    TreeNode* right;
    TreeNode(int v) : val(v), left(nullptr), right(nullptr) {}
};

vector<int> postorderTwoStacks(TreeNode* root) {
    vector<int> out;
    if (root == nullptr) return out;

    stack<TreeNode*> s1, s2;
    s1.push(root);
    while (!s1.empty()) {
        TreeNode* node = s1.top();
        s1.pop();
        s2.push(node);
        if (node->left)  s1.push(node->left);
        if (node->right) s1.push(node->right);
    }
    while (!s2.empty()) {
        out.push_back(s2.top()->val);
        s2.pop();
    }
    return out;
}

int main() {
    TreeNode* root = new TreeNode(1);
    root->left = new TreeNode(2);
    root->right = new TreeNode(3);

    for (int x : postorderTwoStacks(root)) cout << x << " ";
    cout << "\n";
    return 0;
}
\end{lstlisting}

\begin{complexitybox}
\textbf{Time Complexity.} $O(n)$ -- every node passes through both stacks
once. \textbf{Space Complexity.} $O(n)$ -- stack~2 accumulates all $n$
nodes before the final reversal, which is the price of the simpler logic.
\end{complexitybox}

\begin{takeawaybox}
Recognising that postorder is a reversed, child-swapped preorder turns a
fiddly problem into two easy passes. The trade is memory: it holds all $n$
nodes at once. When an interviewer demands $O(h)$ space, use the one-stack
method in the next section instead.
\end{takeawaybox}
